% Root continuation; not input pending independent literal audit.
\subsection{Strong pairs with an adjoined trivial endpoint}

\begin{lemma}\label{C:active-corner-fixed-state}
Let $S$ be a finite strongly connected two-object category with
rectangular-band minimal ideal. Suppose that its row counts are
$(r,1)$, with the larger row set at $a$. If a finite right
$S$-act $F$ has an $\operatorname{End}_S(a)$-fixed state in $F(a)$,
every degree-two rational product of its act cone is zero.
\end{lemma}
\begin{proof}
In the natural categorical balanced-tensor presentation of
$H_2(BS;\mathbb Q)$, the row augmentation at the other object
is zero. Thus every class is represented by a band-grid cycle
at $a$: inclusion of the corner monoid at $a$ surjects onto
$H_2(BS;\mathbb Q)$. The chosen fixed state gives an actual
functor from that corner monoid to the element category of $F$.
Its composite with the projection to $S$ is that corner inclusion.
Consequently $H_2$ of the element category surjects onto $H_2(BS)$.
The cofiber sequence makes every degree-two class of the act cone
relative to $BS$. The relative quotient is a suspension, so its
positive-degree cup products, and their absolute images, vanish.
\end{proof}

\begin{theorem}\label{C:strong-pair-trivial-endpoint21}
A category obtained by adjoining a trivial sink or source to a
strongly connected two-object category cannot have type
$\mathbb CP^2$ with at most twenty-one arrows.
\end{theorem}
\begin{proof}
Take opposites to make the new object a sink, and call the strong
source $S$, its arrow count $s$, and its right act $F$. The two
stalks of $F$ are nonempty. A retract object of $S$ can be deleted
from the whole category, giving an ordinary act cone with fewer
than twenty-two arrows, which is impossible. Assume no such
retract exists.

Proposition~\ref{C:proper-strong-pair-at-least-seventeen} treats
$s\geq17$, and Proposition~\ref{C:strong-pair-fourteen-source-cone}
treats $s\leq14$. Hence $s=15$ or $16$.
If $H^2(BS;\mathbb Q)=0$, the relative-suspension argument
already gives zero products. Otherwise both total minimal row
and column counts are at least three. A nontrivial minimal group
would require at least twenty arrows including the two identities.
Thus the group is trivial. If every object has at most two rows,
Theorem~\ref{C:categorical-two-row-act-cups} applies.

The only remaining shape is a rectangular band with total row
count four and column count three, with row counts $(3,1)$.
Indeed a larger total row count would require at least
$5\cdot3+2=17$ arrows, while total column count at least four
would require at least $4\cdot4+2=18$. Write $a$ for the
three-row object and $b$ for the one-row object. Nonzero second
homology forces the column counts to be $(2,1)$ at $(a,b)$:
in the other distribution each objectwise row-column augmentation
tensor is zero. The minimal band and identities use fourteen
arrows. There are at most two extra internal arrows.

The corner at $b$ has a minimal two-sided zero $e_b$.
If every incoming arrow $a\to b$ belongs to the minimal ideal,
it is fixed by $e_b$, and
Lemma~\ref{C:one-sided-corner-deletion} deletes $b$ from the whole
category. The result is an ordinary act cone with fewer than
twenty-two arrows, again impossible.
We may therefore assume that at least one extra arrow goes
from $a$ to $b$.

If there is no extra endomorphism at $a$, its endomorphism monoid
is precisely the identity-adjoined $3\times2$ band. The left row
augmentation module $V$ in the projective sequence
$0\to V\to P\to\mathbf1\to0$ is supported only at $a$.
It is two-dimensional there, every nonidentity corner element
acts by zero, and every cross map is zero. Consequently
$V=L\oplus L$ for a one-dimensional rational $S$-module $L$.
Using the natural act-cone Tor formula and dimension shifting,
\[
 H_4(B\mathcal C(S,F);\mathbb Q)
   \cong\operatorname{Tor}^{\mathbb QS}_3(L_F,\mathbf1)
   \cong\operatorname{Tor}^{\mathbb QS}_2(L_F,V)
   \cong\operatorname{Tor}^{\mathbb QS}_2(L_F,L)^{\oplus2}.
\]
Its dimension is even, contradicting the fourth Betti number one
of a projective plane. This covers every remaining $s=15$ case.

Finally suppose an extra endomorphism at $a$ exists. Along with
the extra incoming arrow at $b$ it uses both extra arrows, so
$s=16$. Every internal arrow $b\to a$ is minimal and is fixed
on the right by $e_b$, because the column set at $b$ is a singleton.
The total state count satisfies $|F(a)|+|F(b)|\leq4$.
If $|F(b)|=1$, its external outgoing arrow is also fixed by
$e_b$, and the outgoing version of the one-sided deletion lemma
deletes $b$. If $|F(a)|=1$, apply
Lemma~\ref{C:active-corner-fixed-state}.
The only remaining positive pair is $(|F(a)|,|F(b)|)=(2,2)$,
excluded by Lemma~\ref{C:two-states-objectwise-act-cups}.
Every case is impossible.
\end{proof}
